% ============================================================
% Section 45: 一维单原子链的长波极限与连续弹性波方程
% ============================================================
\newpage
\section{固体理论：一维单原子链的长波极限——从离散晶格到连续弹性波}
\label{sec:long-wave-continuum}

\begin{modelbox}[题目：证明长波近似下离散晶格退化为弹性波方程]
证明在长波近似下，一维单原子链的运动方程可化成弹性波方程
\begin{equation}
\frac{\partial^2 u}{\partial t^2} = v^2\frac{\partial^2 u}{\partial x^2}.
\end{equation}
\end{modelbox}

\begin{tipbox}[考点快速分析]
\begin{itemize}
    \item \textbf{多近邻相互作用}：不只考虑最近邻，一般性地包含第 $p$ 近邻的力常数 $\beta_p$
    \item \textbf{对称性}：$\beta_{-p}=\beta_p$（作用与反作用）
    \item \textbf{长波极限}：$qa\ll1$，离散结构被``抹平''，晶体表现为连续介质
    \item \textbf{核心结论}：声学支在长波极限下就是弹性波，$\omega=vq$
\end{itemize}
\end{tipbox}

\begin{derivationbox}[标准解题过程]
\textbf{Step 1: 离散运动方程}

设晶格常数为 $a$，原子质量为 $m$。第 $n$ 个原子与第 $n+p$ 个原子间的恢复力常数为 $\beta_p$（$p=\pm1,\pm2,\dots$），且由对称性 $\beta_{-p}=\beta_p$。第 $n$ 个原子的运动方程为
\begin{equation}
m\ddot u_n = \sum_{p\neq0}\beta_p(u_{n+p}-u_n).
\label{eq:general-motion}
\end{equation}

\textbf{Step 2: 格波解与色散关系}

设格波试探解 $u_n=u_0\,e^{i(\omega t-nqa)}$，则 $u_{n+p}=u_n\,e^{-ipqa}$。代入式 \eqref{eq:general-motion} 并消去公共因子：
\begin{equation}
-m\omega^2 = \sum_{p\neq0}\beta_p\bigl(e^{-ipqa}-1\bigr).
\end{equation}

将求和按 $p>0$ 与 $p<0$ 配对，利用 $\beta_{-p}=\beta_p$：
\begin{align}
-m\omega^2 &= \sum_{p>0}\Bigl[\beta_p\bigl(e^{-ipqa}-1\bigr)+\beta_p\bigl(e^{ipqa}-1\bigr)\Bigr] \nonumber\\
&= \sum_{p>0}\beta_p\bigl(e^{ipqa}+e^{-ipqa}-2\bigr) \nonumber\\
&= -2\sum_{p>0}\beta_p\bigl(1-\cos pqa\bigr) \nonumber\\
&= -4\sum_{p>0}\beta_p\sin^2\frac{pqa}{2}.
\end{align}
故色散关系为
\begin{equation}
\boxed{\omega^2 = \frac{4}{m}\sum_{p>0}\beta_p\sin^2\frac{pqa}{2}.}
\label{eq:general-dispersion}
\end{equation}

\textbf{Step 3: 长波极限 $qa\ll1$}

当 $qa\ll1$ 时，$\sin(pqa/2)\approx pqa/2$，代入式 \eqref{eq:general-dispersion}：
\begin{equation}
\omega^2 \approx \frac{4}{m}\sum_{p>0}\beta_p\Bigl(\frac{pqa}{2}\Bigr)^2 = \Bigl(\frac{a^2}{m}\sum_{p>0}\beta_p p^2\Bigr)q^2.
\end{equation}
定义弹性波速
\begin{equation}
\boxed{v^2 = \frac{a^2}{m}\sum_{p>0}\beta_p p^2,}
\label{eq:elastic-velocity}
\end{equation}
则长波色散简化为
\begin{equation}
\boxed{\omega^2 = v^2 q^2 \quad\text{或}\quad \omega = vq.}
\label{eq:linear-dispersion}
\end{equation}
这正是弹性波的线性色散：相速等于群速，均为常数 $v$。

\textit{最近邻特例}：若只计 $p=1$，$\beta_1=\beta$，则 $v=a\sqrt{\beta/m}$，与一维单原子链的熟知结果一致。

\textbf{Step 4: 过渡到连续介质方程}

长波条件下波长 $\lambda=2\pi/q\gg a$，离散位置 $x_n=na$ 可近似为连续变量 $x$。将格波写为连续形式
\begin{equation}
u(x,t) = u_0\,e^{i(\omega t-qx)},
\end{equation}
分别求二阶偏导：
\begin{equation}
\frac{\partial^2 u}{\partial t^2} = -\omega^2 u, \qquad \frac{\partial^2 u}{\partial x^2} = -q^2 u.
\end{equation}
将长波色散 $\omega^2=v^2q^2$ 代入时间导数：
\begin{equation}
\frac{\partial^2 u}{\partial t^2} = -v^2q^2 u = v^2\frac{\partial^2 u}{\partial x^2}.
\end{equation}
即得到\textbf{一维弹性波方程}
\begin{equation}
\boxed{\frac{\partial^2 u}{\partial t^2} = v^2\frac{\partial^2 u}{\partial x^2}.}
\end{equation}
\end{derivationbox}

\begin{formulabox}[核心公式链]
\begin{center}
\begin{tabular}{cl}
\toprule
步骤 & 结果 \\
\midrule
离散运动方程 & $\displaystyle m\ddot u_n = \sum_{p\neq0}\beta_p(u_{n+p}-u_n)$ \\[8pt]
色散关系 & $\displaystyle \omega^2 = \frac{4}{m}\sum_{p>0}\beta_p\sin^2\frac{pqa}{2}$ \\[8pt]
长波极限 & $\displaystyle \omega^2 = v^2 q^2$，$\displaystyle v^2=\frac{a^2}{m}\sum_{p>0}\beta_p p^2$ \\[8pt]
连续化 & $\displaystyle \partial_t^2 u = v^2\partial_x^2 u$ \\
\bottomrule
\end{tabular}
\end{center}
\end{formulabox}

\begin{insightbox}[物理图像与概念深化]
\textbf{1. 为什么长波极限下离散退化为连续？}

当波长 $\lambda\gg a$ 时，波在一个周期内跨越了成千上万个原子。此时观察者``看不清''单个原子的位移，只能感知到一段原子团的\textbf{集体平滑运动}。离散的结构细节被平均掉，晶格表现得像一根均匀的弹性弦。

\textbf{2. 弹性波速 $v$ 的微观表达式}

式 \eqref{eq:elastic-velocity} 表明：宏观弹性波速由微观参数共同决定：
\begin{itemize}
    \item $a$：晶格常数（原子间距）
    \item $m$：原子质量（惯性）
    \item $\beta_p$：力常数（刚度）
    \item $p^2$：权重因子（更远邻居的贡献按距离平方放大）
\end{itemize}
值得注意的是，更远邻居（大 $p$）虽然力常数 $\beta_p$ 通常衰减很快，但权重 $p^2$ 增长，因此在某些长程相互作用体系中，次近邻、第三近邻的贡献不可忽略。

\textbf{3. 与光学支的本质区别}

声学支在长波极限下 $\omega\to0$，对应原胞质心的整体平移，因此能退化为连续介质的弹性波。光学支在长波极限下 $\omega\neq0$，对应原胞内原子的相对振动，这种\textbf{内自由度}在连续介质模型中没有对应物——连续介质假设的前提是``物质不可区分''，而光学支恰恰依赖于物质内部的相对运动。

\textbf{4. 高阶修正：从弹性波到色散波}

若保留 $qa$ 的高阶项：
\begin{equation}
\sin^2\frac{pqa}{2} = \Bigl(\frac{pqa}{2}\Bigr)^2 - \frac{1}{3}\Bigl(\frac{pqa}{2}\Bigr)^4 + \cdots
\end{equation}
则
\begin{equation}
\omega^2 = v^2 q^2 - \gamma q^4 + \cdots, \qquad \gamma = \frac{a^4}{12m}\sum_{p>0}\beta_p p^4.
\end{equation}
$q^4$ 项导致波的\textbf{色散}（不同频率分量传播速度不同），这是离散晶格偏离连续弹性介质的标志。在 X 射线声学、超声波等领域，这种色散效应可被观测。
\end{insightbox}

\begin{thinkbox}[考场速记：长波极限问题的证明模板]
遇到``证明长波极限下退化为弹性波/连续方程''类问题：
\begin{enumerate}
    \item \textbf{写离散方程}：$m\ddot u_n = \sum_p \beta_p(u_{n+p}-u_n)$
    \item \textbf{设格波解}：$u_n = u_0 e^{i(\omega t-nqa)}$
    \item \textbf{求色散}：配对正负 $p$，利用 $1-\cos\theta=2\sin^2(\theta/2)$
    \item \textbf{长波展开}：$\sin(pqa/2)\approx pqa/2$，得到 $\omega=vq$
    \item \textbf{连续化}：$x=na$，$u(x,t)=u_0 e^{i(\omega t-qx)}$，代入 $\omega^2=v^2q^2$ 即得波动方程
\end{enumerate}
\textit{检验}：若只计最近邻，应能退化为 $v=a\sqrt{\beta/m}$，与一维单原子链熟知结果一致。
\end{thinkbox}
